\section{Introduction}

This report is about one graphics card and what it can actually do. The card is
an RTX 5070, a Blackwell part with 12\,GB of GDDR7 that also happens to be
driving the display on the machine I built this on. The vendor quotes a peak
around \SI{31}{TFLOPS} of FP32 and a memory bandwidth around \SI{672}{GB/s}. Those
two numbers are the headline of any spec sheet, and on their own they say almost
nothing about how fast a given kernel will run, because most kernels are limited
not by raw compute but by how much data they have to move to feed that compute.

The roofline model exists to make that trade off visible. It plots attainable
performance against arithmetic intensity, the ratio of floating point work to
bytes moved, and bounds every kernel under two ceilings at once: a flat compute
ceiling and a sloped memory ceiling. Where a kernel sits under that roof tells
you which resource it is starved on, and therefore what would make it faster.

The contribution here is not another fast GEMM. It is a ladder of kernels that
walks the arithmetic intensity axis on purpose, from a SAXPY that touches memory
and does almost no arithmetic, through a transpose that does no floating point
work at all, up to a register blocked GEMM that reuses each loaded value many
times. For every rung I measure where it lands, and I explain that position with
the specific hardware counters behind it, not with a generic statement about
what tiling usually does. The two ceilings I draw are both real: one derived
from the device properties, one measured by pushing the card as hard as I can
with a microbenchmark, and I am honest about the distance between them.
